% ECONOMICS_PROBLEM: {"id": "F16-277", "title": "Compensation for trade adjustment", "jel_code": "F16", "source_id": "OP-0341"}
% !TeX program = xelatex
\documentclass[11pt,a4paper]{article}
\usepackage[margin=23mm]{geometry}
\usepackage{fontspec}
\setmainfont{Latin Modern Roman}
\usepackage{amsmath,amssymb,mathtools}
\usepackage{xcolor,enumitem,booktabs,longtable,array,xurl,hyperref}
\definecolor{muted}{HTML}{53636C}
\hypersetup{colorlinks=true,urlcolor=blue,linkcolor=blue}
\setlength{\parindent}{0pt}
\setlength{\parskip}{5pt}
\newcommand{\R}{\mathbb R}
\newcommand{\E}{\mathbb E}
\newcommand{\N}{\mathbb N}
\newcommand{\ind}{\mathbf 1}
\newcommand{\Var}{\operatorname{Var}}
\newcommand{\Cov}{\operatorname{Cov}}
\newcommand{\supp}{\operatorname{supp}}
\DeclareMathOperator*{\argmin}{arg\,min}
\DeclareMathOperator*{\argmax}{arg\,max}
\newcommand{\entrymeta}[3]{{\sffamily\small\color{muted}\textbf{\texttt{#1}}\quad|\quad #2\par #3\par}}
\newcommand{\Problemref}[1]{\texttt{#1}}
\newcommand{\JELref}[2]{\texttt{#2} (JEL \texttt{#1})}
\begin{document}
\section*{Compensation for trade adjustment}
\textbf{Problem ID:} \texttt{F16-277}\quad \textbf{JEL:} \texttt{F16}\par
% BEGIN ECONOMICS STATEMENT
\label{problem:OP-0341}
{\sffamily\small\color{muted}Original subject group: Trade, globalization and geoeconomics\par}
\label{definitions:problem:30007623}
\textbf{Audit classification:} Background; proposed extension. The citation mixed the 2021 NBER WP with its 2024 journal version. Both are separately corrected. The source discusses compensation difficulties; it does not verify this exact optimal multicomponent program as an open theorem.
\subsection*{Definitions and precise research target}
Fix an externally specified trade shock affecting workers and locations. A feasible adjustment policy \(a\) allocates income transfers, retraining, mobility support and place-based investment under budget \(B\). Workers choose participation, occupations and residence; firms adjust entry and vacancies. Let \(V_i(a)\) be worker \(i\)'s discounted expected utility after these responses. Fix welfare weights \(\omega_i\geq0\), a discount convention and baseline \(a=0\). The design task is
\[A_B=\{a\in A:C(a)\leq B\},\qquad W_*=\sup_{a\in A_B}\sum_i\omega_iV_i(a),\]
where \(C(a)\) measures fiscal outlays and administration under the declared budget convention, with financing distortions separately and consistently included in welfare and \(A\) incorporates legal and delivery constraints. Identify or bound the causal effects of each instrument on consumption, earnings, employment and movement, including spillovers on untreated workers and places. Distinguish compensating loss from changing future opportunity: equal transfers can produce different long-run adjustments. Assess complementarity between support components and heterogeneous take-up. State which estimates transport from observed programs to a larger policy scale and which rely on equilibrium assumptions. A policy helping the treated group can still reduce weighted welfare through financing or displacement. No claim is made that a universally optimal compensation mechanism is already a canonical unresolved theorem.

\textbf{Additional formal definitions and scope (editorial).} Take \(A\) to consist of measurable transfer, training, mobility and investment rules conditional on the information the government actually observes. State whether losses and preferences are observable or only reported. Worker \(V_i\) includes consumption, labor and moving costs with cardinal normalization fixed before weighting. Define \(C(a)\) as the fiscal outlay plus administration relevant to the budget; distortionary-financing deadweight loss belongs in welfare or a separate cost constraint, but not both without a stated accounting convention. Tax receipts and transfers cancel as transfers in aggregate resource accounting. If \(A\) is compact and the weighted value is upper semicontinuous, the optimum exists by standard compactness; the unresolved empirical design issues concern response, eligibility, leakage, spillovers and financing. Private information adds a distinct incentive-compatible mechanism-design problem. For a nonempty feasible set and a finite optimal value, the design target is an \(\varepsilon\)-optimal feasible rule for each specified \(\varepsilon>0\): its cost is at most the infimum plus \(\varepsilon\), or its welfare is at least the supremum minus \(\varepsilon\). Report infeasibility or an unbounded value separately. An exact optimizer is requested only under additional attainment assumptions; it need not exist for the general rule class.
\subsection*{Variants and assumption changes}
\begin{enumerate}
\item \textbf{Observable losses and transfers (editorial).} Observe worker types, restrict support to lump-sum transfers and include distortionary financing; determine whether feasible compensation benefits the weighted population.
\item \textbf{Unobserved losses and adjustment programs (editorial).} Add private information, training and mobility; distinguish a causal program-evaluation problem from an incentive-compatible policy-design extension.
\end{enumerate}
\subsection*{References and source audit}
\begin{enumerate}
\item NBER verifies Rodrik, 2021 WP 29507 and DOI; it is distinct from the 2024 journal publication originally linked. Locator: 2021 NBER abstract; author-hosted version, compensation discussion printed pp.3--4. Support: Trade-distribution and compensation background. NBER indexed metadata and author-hosted text inspected. URL: \url{https://www.nber.org/papers/w29507}.
\item Publisher verifies Rodrik, 2024, Oxford Open Economics 3(Supplement 1):i1076--i1082, DOI 10.1093/ooec/odad048; not the 2021 NBER DOI. Locator: Compensation discussion; author-hosted version printed pp.3--4. Support: Compensation background; no exact optimal-program open theorem. Publisher indexed text and author-hosted version inspected; repeated publisher fetch failed. URL: \url{https://academic.oup.com/ooec/article/3/Supplement_1/i1076/7708093}.
\end{enumerate}
\begin{enumerate}\raggedright
\item\hypertarget{definitions:source:30007623:1}{} Dani Rodrik. 2021. \emph{A Primer on Trade and Inequality}. 2021 NBER abstract; author-hosted version, compensation discussion printed pp.3--4.
\url{https://www.nber.org/papers/w29507}.
DOI: \url{https://doi.org/10.3386/w29507}.
\item\hypertarget{definitions:source:30007623:2}{} Dani Rodrik. 2024. \emph{A primer on trade and inequality}. Compensation discussion; author-hosted version printed pp.3--4.
\url{https://academic.oup.com/ooec/article/3/Supplement_1/i1076/7708093}.
DOI: \url{https://doi.org/10.1093/ooec/odad048}.
\end{enumerate}

\subsection*{Common conventions: Source mathematical conventions}

These conventions apply to each entry unless explicitly replaced.

\textbf{Models and identification.} A primitive model \(m\) specifies all probability laws, feasible choices, information, equilibrium and measurement rules used by its target. The admissible class \(\mathcal M\) is fixed independently of outcomes. If \(P_O(m)\) is its observable probability law and \(\tau(m)\) the target, its identified set at observable law \(P\) is
\[
 \mathcal I_\tau(P)=\{\tau(m):m\in\mathcal M,\ P_O(m)=P\}.
\]
Point identification means that this set is a singleton. An empty set signals incompatibility of data and assumptions. Coordinate bounds need not describe the joint set. Bounds are sharp only if every retained target is attainable or arbitrarily approachable by compatible models. Finite bounds require explicit restrictions on unobserved quantities; unrestricted confounding, hidden wealth, extrapolation or arbitrary equilibrium selection can yield uninformative sets.

\textbf{Causal interventions.} A potential outcome \(Y_i(p)\) is the outcome under a complete policy regime \(p\), including timing, financing, information and market spillovers. The target population and units are fixed before treatment. With interference, use the complete assignment vector or a justified exposure mapping; individual no-interference notation alone cannot identify an economy-wide intervention. A difference of mean potential outcomes does not require the cross-world joint distribution. Conditioning on post-treatment migration, survival, enrollment or employment changes the estimand and needs separate justification.

\textbf{Statistical statements.} An estimator, confidence set, forecast or test specifies a sampling law, model class and loss/coverage criterion. Uniform \((1-\alpha)\)-coverage means \(\inf_{m\in\mathcal M}\Pr_m\{\tau(m)\in C_n\}\geq1-\alpha\), or an explicitly asymptotic version. The whole parameter space trivially has coverage, so informativeness requires a stated width, risk or power objective. A forecasting improvement is relative to a named benchmark, information set and evaluation population.

\textbf{Equilibrium and optimization.} An equilibrium comprises feasible plans, optimality, beliefs where needed, budget constraints and all market/resource clearing. The primitives or a stated benchmark define each correspondence \(\mathcal E(p;m)\). Nonemptiness is assumed or proved, never inferred just from compact policy sets. A selected-equilibrium target uses a declared selection rule; a robust target quantifies over all permitted equilibria. Use suprema or infima unless attainment is established. An \(\varepsilon\)-optimal policy is feasible and within \(\varepsilon\) of the finite optimum value. Report an empty feasible set separately.

\textbf{Welfare and accounting.} Normative utility, interpersonal normalization, social weights, numeraire, discounting and the use of public receipts are primitives. Behavioral utility need not coincide with the welfare criterion. Household welfare includes ownership income and rents once; adding profits again to compensating variation generally double-counts them. A monetary equivalent is defined by an explicit transfer and price convention, with monotonicity and range conditions. Average effects, mechanism contributions, transfers and welfare losses are different targets. Infinite sums require convergence and dynamic feasible plans require terminal or no-Ponzi conditions.

\textbf{Source versions.} A working-paper date, revision date and journal publication year are distinguished. A DOI for a journal version is not automatically the DOI of an earlier manuscript. Locators refer to the stated version and printed pages where specified. A metadata-only check does not verify a claim about a full-text conclusion. Every entry's source subsection narrows the provenance claim accordingly.
% END ECONOMICS STATEMENT
\end{document}
